% Solutions of 11-induction-and-termination.
\section*{Chapter \ref{ch:b1:induction-and-termination} --- Induction, Recursion and Termination}

\begin{solution}{exo:b1:induction-and-termination:1}
Empty tree: \omterm{def:b1:induction-and-termination:height}{size} 0, \omterm{def:b1:induction-and-termination:height}{height} $-1$. Single node: 1 and 0. A node with a left child only: 2 and 1. The tree \code{t}: 4 nodes
(5, 2, 9, 7), and the longest path 5, 9, 7 has two links: \omterm{def:b1:induction-and-termination:height}{height} 2, as \cref{out:b1:induction-and-termination:inductive}
prints.
\end{solution}

\begin{solution}{exo:b1:induction-and-termination:2}
$(252, 105) \to (105, 42) \to (42, 21) \to (21, 0)$: the \omterm{def:b1:induction-and-termination:variant}{variant} $b$ takes the values 105, 42, 21 and 0, and the
result is 21.
\end{solution}

\begin{solution}{exo:b1:induction-and-termination:3}
\omoutput[2-3]{code/bachelor-1/11-induction-and-termination/outputs/exo_lcrs.txt}{The two \omterm{def:b1:induction-and-termination:binary-tree}{binary trees}.}
In the first tree, 2 is the first child of 1 (left of 1), 4 the first child of 2 (left of 2), and 3 the next sibling of
2 (right of 2). In the second, the three \omterm{def:b1:induction-and-termination:rooted-tree}{children} of 1 form a chain of right links under the left child 2.
\end{solution}

\begin{solution}{exo:b1:induction-and-termination:4}
(a) \omterm{def:b1:induction-and-termination:well-founded}{Well-founded}: a descending chain of integers greater than $-10$ starting at $x_0$ has at most $x_0 + 10$ steps.
(b) \omterm{def:b1:induction-and-termination:well-founded}{Well-founded}: a proper divisor of a positive integer is smaller than it, so a descending chain is a descending chain
in $\mathbb{N}$. (c) \omterm{def:b1:induction-and-termination:well-founded}{Well-founded}: a proper prefix is shorter. (d) Not \omterm{def:b1:induction-and-termination:well-founded}{well-founded}: $1 > 1/2 > 1/4 > \cdots$.
\end{solution}

\begin{solution}{exo:b1:induction-and-termination:5}
\omoutput[2-7]{code/bachelor-1/11-induction-and-termination/outputs/exo_variant.txt}{The calls.}
The call \code{steps(n, 1)} made from \code{steps(n, 0)} keeps $n$: $n$ is not a \omterm{def:b1:induction-and-termination:variant}{variant}. The pair $(n, 1 -
\mathit{step})$ is one, in the \omterm{def:b1:induction-and-termination:lex}{lexicographic order} of $\mathbb{N}^2$: from $(n, 1)$ to $(n, 0)$ on the first kind of
call, from $(n, 0)$ to $(n - 1, 0)$ on the second. By induction along this order, \code{steps(n, step)} returns $n$ for
every natural number $n$.
\end{solution}

\begin{solution}{exo:b1:induction-and-termination:6}
\omoutput[2-4]{code/bachelor-1/11-induction-and-termination/outputs/exo_height0.txt}{The two conventions compared.}
With that convention every \omterm{def:b1:induction-and-termination:height}{height} is one more, $h' = h + 1$. The bound $n \leq 2^{h + 1} - 1$ still holds but is no
longer reached; $n \geq h + 1$ fails, for the empty tree and for a single node. Rewritten: $h' \leq n \leq 2^{h'} - 1$,
and at most $2^{h' - 1}$ leaves (\cref{exo:b1:induction-and-termination:9}).
\end{solution}

\begin{solution}{exo:b1:induction-and-termination:7}
\omcode{code/bachelor-1/11-induction-and-termination/exercises/count_paths.py}{4}{9}{Paths from the \omterm{def:b1:induction-and-termination:rooted-tree}{root} to a \omterm{def:b1:induction-and-termination:rooted-tree}{leaf}}
By \omterm{def:b1:induction-and-termination:structural}{structural induction}: the empty tree has no path and no \omterm{def:b1:induction-and-termination:rooted-tree}{leaf}; a single node has one path and one \omterm{def:b1:induction-and-termination:rooted-tree}{leaf}; otherwise a
path from the \omterm{def:b1:induction-and-termination:rooted-tree}{root} to a \omterm{def:b1:induction-and-termination:rooted-tree}{leaf} goes into exactly one of the two \omterm{def:b1:induction-and-termination:rooted-tree}{subtrees}, and continues as a path from the \omterm{def:b1:induction-and-termination:rooted-tree}{root} of that
\omterm{def:b1:induction-and-termination:rooted-tree}{subtree} to a \omterm{def:b1:induction-and-termination:rooted-tree}{leaf} of it, so the counts add, as the leaves do. On a \omterm{def:b1:induction-and-termination:binary-tree}{full} tree with $i$ \omterm{def:b1:induction-and-termination:rooted-tree}{internal nodes} there are $i + 1$
leaves (\cref{prop:b1:induction-and-termination:bounds}), hence $i + 1$ paths. The test \code{tests/test_trees.py}
checks both on random trees.
\end{solution}

\begin{solution}{exo:b1:induction-and-termination:8}
Empty tree: $\text{len}([\,]) = 0 = \text{size}$. Node $(\ell, k, r)$: $\text{len}(\mathit{flatten}(\ell) + [k] +
\mathit{flatten}(r)) = \text{len}(\mathit{flatten}(\ell)) + 1 + \text{len}(\mathit{flatten}(r)) = \text{size}(\ell) + 1 +
\text{size}(r)$ by the induction hypothesis, which is the \omterm{def:b1:induction-and-termination:height}{size} of the node.
\end{solution}

\begin{solution}{exo:b1:induction-and-termination:9}
By \omterm{def:b1:induction-and-termination:structural}{structural induction} on non-empty trees. A single node: \omterm{def:b1:induction-and-termination:height}{height} 0, one \omterm{def:b1:induction-and-termination:rooted-tree}{leaf}, $2^0 = 1$. A node with one empty
\omterm{def:b1:induction-and-termination:rooted-tree}{subtree}: its leaves are those of the other \omterm{def:b1:induction-and-termination:rooted-tree}{subtree}, of \omterm{def:b1:induction-and-termination:height}{height} $h - 1$, so at most $2^{h - 1} \leq 2^h$. A node with two
non-empty \omterm{def:b1:induction-and-termination:rooted-tree}{subtrees} of \omterm{def:b1:induction-and-termination:height}{heights} at most $h - 1$: at most $2^{h - 1} + 2^{h - 1} = 2^h$ leaves. The bound is reached exactly
by the \emph{perfect} trees, in which every \omterm{def:b1:induction-and-termination:rooted-tree}{internal node} has two \omterm{def:b1:induction-and-termination:rooted-tree}{children} and every \omterm{def:b1:induction-and-termination:rooted-tree}{leaf} has \omterm{def:b1:induction-and-termination:height}{depth} $h$.
\end{solution}

\begin{solution}{exo:b1:induction-and-termination:10}
In a descending chain for shortlex, the lengths never increase, so from some point on they are constant, equal to some
$L$ (lengths are natural numbers). From then on the chain is strictly descending among the words of length $L$, of
which there are finitely many: it is finite. For the dictionary order, $b > ab > aab > aaab > \cdots$: at each step
the first letters compare $a < b$ at the first difference, and the chain never ends.
\end{solution}

\begin{solution}{exo:b1:induction-and-termination:11}
\omcode{code/bachelor-1/11-induction-and-termination/exercises/sort_game.py}{1}{20}{The sorting game}
A swap of an adjacent inverted pair $x_i > x_{i+1}$ removes that inversion; every other pair keeps its two elements, in
the same order, so the number of inversions decreases by exactly one. It is a \omterm{def:b1:induction-and-termination:variant}{variant} in $\mathbb{N}$, so every game
ends. It ends when no adjacent pair is inverted, that is when the list is sorted; and it takes exactly as many moves as
the starting list has inversions, whatever the choices. The test plays with the leftmost and the rightmost choice and
checks the result and the count on random lists.
\end{solution}

\begin{solution}{exo:b1:induction-and-termination:12}
Use \omterm{def:b1:induction-and-termination:well-founded}{well-founded} induction on $m = 101 - n \geq 1$ for $n \leq 100$, proving that $f(n)$ \omterm{def:b1:induction-and-termination:variant}{terminates} with value 91.
First, if $90 \leq n \leq 100$, then $n + 11 > 100$, so $f(n + 11) = n + 1$ and $f(n) = f(n + 1)$; with $f(101) = 91$,
going down from 100 to 90 gives $f(n) = 91$. If $n < 90$, then $n + 11 \leq 100$ and $101 - (n + 11) < m$, so by
induction $f(n + 11)$ \omterm{def:b1:induction-and-termination:variant}{terminates} with 91, and $f(n) = f(91) = 91$ by the first case. For $n > 100$ the function
returns at once. The test checks the values from $-50$ to 200.
\end{solution}

\begin{solution}{pb:b1:induction-and-termination:1}
\textbf{1.} The base rule: a position with no entry above 0 is a play by itself. The construction rule: if $s'$ is one
move away from $s$ and $p$ is a play from $s'$, then $s$ followed by $p$ is a play from $s$.

\textbf{2.} From $(2)$: $(2), (1, 1), (0, 0, 1), (0, 0, 0, 0)$ and $(2), (1, 1), (1, 0, 0), (0, 0, 0, 0)$. From $(1, 1)$: the
same two without their first position.

\textbf{3.} The only move from $(4)$ leads to $(3, 3)$, so every play from $(4)$ is $(4)$ followed by a play from
$(3, 3)$.

\textbf{4.} It follows the inductive definition of plays: a final position has one play, of length 0; otherwise the
lengths are one more than those of the plays from each position one move away.

\textbf{5.} The same position is reached by many orders of moves; without the cache, \code{plays} would recompute it
each time, and the count would cost as many calls as there are plays, 21\,964\,800. With the cache, each position is
computed once.

\textbf{6.} No: exchanging entries exchanges the moves, and gives a one-to-one correspondence between the plays of the
two positions, which preserves lengths. The number of plays is the same too.

\textbf{7.} The sum changes by $2(k - 1) - k = k - 2$, which is positive for $k \geq 3$: $(3) \to (2, 2)$ raises it from
3 to 4. The number of entries grows by one at each move. The largest entry can stay the same: $(2, 2) \to (1, 1, 2)$.

\textbf{8.} A move adds only entries $k - 1$, smaller than an existing entry.

\textbf{9.} The move lowers $c_k$ by one and raises $c_{k - 1}$ by two, leaving the others. The first component where
the vectors differ is $c_k$, which decreased.

\textbf{10.} The \omterm{def:b1:induction-and-termination:lex}{lexicographic order} on $\mathbb{N}^{K + 1}$ is \omterm{def:b1:induction-and-termination:well-founded}{well-founded}, and a play gives a descending chain of
vectors $c(s)$: it is finite.

\textbf{11.} No: $c_{k - 1}$ increases, so the new vector is not below the old in the \omterm{def:b1:induction-and-termination:lex}{product order}. Only the existence
of a \omterm{def:b1:induction-and-termination:well-founded}{well-founded order} in which $c$ decreases matters.

\textbf{12.} Part II proves that every play is finite but gives no bound: a lexicographic chain can be as long as one
likes. The table says that every play from $(3, 1, 0)$ has 8 moves.

\textbf{13.} From $(0)$ no move is possible. From $(k)$ the first move gives $(k - 1, k - 1)$, and each copy then needs
$M(k - 1)$ moves of its own: $M(k) = 1 + 2M(k - 1)$.

\textbf{14.} $M(0) = 0 = 2^0 - 1$, and $1 + 2(2^{k - 1} - 1) = 2^k - 1$. The table: 1, 3, 7, 15 for $k = 1, \dots, 4$.

\textbf{15.} Each move acts on one entry, and the entries it creates descend from it. Restricted to the descendants
of entry $i$, the moves of a play form a play from $(a_i)$, and the whole play interleaves these $n$ plays.

\textbf{16.} The moves of the whole play are the moves of the $n$ sub-plays: $\sum_i M(a_i) = \sum_i (2^{a_i} - 1)$.

\textbf{17.} Question 13 assumed that $M(k - 1)$ is well defined. The proof is an induction on $k$: assuming every play
from $(k - 1)$ has $2^{k - 1} - 1$ moves, question 15 applied to $(k - 1, k - 1)$ shows that every play from $(k)$ has
$1 + 2(2^{k - 1} - 1)$ moves.

\textbf{18.} 2, 80 and 8. From $(1, 2)$: the play from $(2)$ has 3 moves, in 2 orders, and the single move of $(1)$ can
come in any of 4 places among them: $4 \times 2 = 8$.

\textbf{19.} $-(2^k - 1) + 2(2^{k - 1} - 1) = -2^k + 1 + 2^k - 2 = -1$.

\textbf{20.} Each term $2^{a_i} - 1$ is at least 0, and 0 only for $a_i = 0$; so $\Phi(s) = 0$ exactly when every entry
is 0, when no move is possible.

\textbf{21.} Each move lowers $\Phi$ by exactly one, and the play ends exactly when $\Phi$ reaches 0: a play from $s$ has
$\Phi(s)$ moves, which also proves \omterm{def:b1:induction-and-termination:variant}{termination} again.

\textbf{22.} Part II gives \omterm{def:b1:induction-and-termination:variant}{termination} and nothing more; $\Phi$ gives both, because it decreases by exactly one. The sum
weighs an entry $k$ by $k$, less than the two entries $k - 1$ can weigh; $\Phi$ weighs it by $2^k - 1$, exactly the cost
of the play it will start, so that the move pays one and leaves the rest to the copies.

\textbf{23.} The check covers 84 positions exhaustively, and would have caught an error in the program, or a wrong
formula; the proof covers every position, of any length and any entries, which no check can enumerate.

\textbf{24.} \emph{Every play from $(a_1, \dots, a_n)$ has the same length $\sum_i (2^{a_i} - 1)$}: each move lowers the
potential $\Phi$ by exactly one, and the play ends exactly when $\Phi$ is 0. The table and the test confirm it on all
the positions they try, from $(1)$ with 1 move to $(4)$ and $(3, 3)$, whose 21\,964\,800 plays all have 15 and 14 moves.
\end{solution}
